\chapter{Eigenvalues and Invariant Subspaces}\label{linear:ch:eigen}

\begin{center}\itshape
I have not thought it necessary to undertake the labour of a formal proof of the theorem in the general case of a
matrix of any degree.
\upshape(Arthur Cayley, \emph{A Memoir on the Theory of Matrices}, 1858)
\end{center}

The planets do not quite follow Kepler's ellipses. Each one pulls on the others, and over centuries the ellipses
slowly turn and stretch. Lagrange and Laplace, in the 1770s and 1780s, studied these \emph{secular} variations (from
the Latin \emph{saeculum}, a century) by linearising the equations of motion, and found themselves solving a polynomial
equation whose roots gave the frequencies of the slow oscillations: the \emph{secular equation}. Whether its roots
were real decided whether the solar system was stable or would fly apart, and Laplace argued that they were. In 1829
Cauchy proved it in general: the secular equation of a symmetric system has only real roots. He was computing the
eigenvalues of a symmetric matrix, decades before anyone used either word.

The words came later and from elsewhere. Sylvester coined ``matrix'' in 1850, a womb from which determinants are born.
Cayley, in 1858, wrote the first memoir on matrices as objects in their own right, and stated that every matrix
satisfies its own characteristic equation; he checked it for $2\times2$ and $3\times3$ matrices and, in the sentence
above, declined to prove the rest. Frobenius proved it in 1878. Meanwhile Weierstrass (1868) and Camille Jordan (1870)
found the normal forms of a single linear map, and in 1874 Jordan and Kronecker quarrelled in the \emph{Comptes rendus}
over whose form was better and what a ``canonical'' form even meant. Today, Jordan's name is on the form and
Weierstrass's theory is its proof; the quarrel is forgotten, but its question, what is the simplest matrix of a given
linear map, is the subject of this chapter.

We work with a single operator $A$ on a finite-dimensional space $V$ and study the subspaces it maps into themselves.
Eigenvectors are the one-dimensional ones; when there are enough of them, $A$ is diagonalisable, as in Tome~0. When
there are not, we prove that $A$ can still be made triangular (Schur), that it satisfies its characteristic polynomial
(Cayley--Hamilton, with a proof valid over every field), that $V$ splits into generalised eigenspaces, and that on each of
them $A$ is a scalar plus a nilpotent operator. The structure of a nilpotent operator is then read off a Young diagram
built from the dimensions of the kernels of its powers, and this gives the Jordan form with a picture instead of a
computation.

\begin{think}
A $4\times4$ complex matrix $N$ satisfies $N^4=0$. Up to similarity (a change of basis), how many such matrices are
there? What if $N^2=0$? Try to guess what single invariant would tell two of them apart. Section~\ref{linear:sec:nilpotent}
answers: the matrices with $N^n=0$ in dimension $n$ correspond to the partitions of $n$, so there are $5$ of them for
$n=4$, and $3$ of those satisfy $N^2=0$.
\end{think}

\section{Invariant subspaces}

Let $A$ be an operator on a vector space $V$ of dimension $n$ over a field $\mathbb K$ (we write $A$ both for the
operator and for its matrix in a chosen basis). A subspace $W$ is \emph{invariant} (or \emph{stable}) if $A(W)\subseteq
W$. Then $A$ restricts to an operator $A|_W$ on $W$, and it induces an operator $\bar A$ on the quotient $V/W$,
$\bar A(v+W)=Av+W$, well defined because $A$ maps $W$ into $W$. In a basis of $V$ that starts with a basis of $W$, the
matrix of $A$ is block upper triangular,
\[
A=\begin{pmatrix}A|_W&*\\0&\bar A\end{pmatrix},
\]
where the lower block is the matrix of $\bar A$ in the basis of $V/W$ formed by the classes of the remaining vectors. In
particular the characteristic polynomial factors:
\begin{equation}\label{linear:eq:charfactor}
\chi_A=\chi_{A|_W}\,\chi_{\bar A},\qquad\text{where }\chi_A(X)=\det(XI-A).
\end{equation}
An \emph{eigenvector} is a non-zero $v$ with $Av=\lambda v$, so that $\mathbb Kv$ is an invariant line; $\lambda$ is an
\emph{eigenvalue}, a root of $\chi_A$. Kernels and images of polynomials in $A$ are invariant, because they commute with
$A$: if $P(A)v=0$ then $P(A)Av=AP(A)v=0$.

\begin{proposition}\label{linear:prop:indep}
Eigenvectors $v_1,\dots,v_k$ for distinct eigenvalues $\lambda_1,\dots,\lambda_k$ are linearly independent. Consequently,
if $\chi_A$ has $n$ distinct roots in $\mathbb K$, then $A$ is diagonalisable.
\end{proposition}
\begin{proof}
By induction on $k$. If $\sum_i c_iv_i=0$, apply $A-\lambda_k$: $\sum_{i<k}c_i(\lambda_i-\lambda_k)v_i=0$, so by
induction $c_i(\lambda_i-\lambda_k)=0$ and $c_i=0$ for $i<k$; then $c_kv_k=0$ gives $c_k=0$.
\end{proof}

\section{Triangularisation}

We say that a polynomial \emph{splits} over $\mathbb K$ if it is a product of polynomials of degree $1$. Over $\C$ every
polynomial splits (the fundamental theorem of algebra, Tome~0, Chapter~3).

\begin{theorem}[triangularisation]\label{linear:thm:schur}
An operator $A$ is triangularisable (it has an upper triangular matrix in some basis) if and only if $\chi_A$ splits.
In particular every complex matrix is similar to an upper triangular one. If moreover $V$ has an inner product (over $\R$
or $\C$, Appendix~D) and $\chi_A$ splits, the basis can be chosen orthonormal (Schur's theorem).
\end{theorem}
\begin{proof}
If $A$ is triangular with diagonal $a_{11},\dots,a_{nn}$, then $\chi_A=\prod(X-a_{ii})$ splits. Conversely, by induction on
$n$: if $\chi_A$ splits, it has a root $\lambda$, with an eigenvector $v_1$. Let $W=\mathbb Kv_1$. By~\eqref{linear:eq:charfactor},
$\chi_{\bar A}=\chi_A/(X-\lambda)$ splits, so by induction $V/W$ has a basis $\bar v_2,\dots,\bar v_n$ in which $\bar A$ is
upper triangular. Lift the classes to vectors $v_2,\dots,v_n$. Then $v_1,\dots,v_n$ is a basis of $V$ (Theorem~\ref{linear:thm:quotient}),
and $Av_j\in a_{jj}v_j+\dots+a_{2j}v_2+\mathbb Kv_1$ for $j\ge2$, so the matrix is upper triangular.

For the orthonormal version, take $v_1$ of norm $1$ and replace the quotient by the orthogonal complement $W^\perp$, which
is a complement of $W$ (Appendix~D): let $\pi$ be the orthogonal projection onto $W^\perp$ and apply the induction to
the operator $\pi A|_{W^\perp}$, whose characteristic polynomial is again $\chi_A/(X-\lambda)$ by the same block computation.
An orthonormal basis $v_2,\dots,v_n$ of $W^\perp$ triangularising it, together with $v_1$, triangularises $A$.
\end{proof}

Triangularisation already gives a lot. The diagonal of a triangular matrix of $A$ lists the eigenvalues with their
multiplicities as roots of $\chi_A$ (the \emph{algebraic multiplicities}), so $\operatorname{tr}A=\sum\lambda_i$ and
$\det A=\prod\lambda_i$; and the eigenvalues of $P(A)$ are the $P(\lambda_i)$, because $P(T)$ is triangular with diagonal
$P(t_{ii})$ when $T$ is. Over $\R$, not every operator is triangularisable: a rotation of the plane by $90^\circ$ has no
invariant line at all.

\section{The Cayley--Hamilton theorem}

\begin{theorem}[Cayley--Hamilton]\label{linear:thm:cayleyhamilton}
For every $A\in M_n(R)$, with $R$ any commutative ring (in particular any field), $\chi_A(A)=0$.
\end{theorem}
\begin{proof}
Let $B(X)=\operatorname{adj}(XI-A)$, the adjugate of $XI-A$; its entries are $(n-1)\times(n-1)$ minors of $XI-A$, so
polynomials of degree at most $n-1$ in $X$, and we can write $B(X)=\sum_{k=0}^{n-1}B_kX^k$ with matrices $B_k\in
M_n(R)$. The adjugate identity (Corollary~\ref{linear:cor:adjugate}, proved there over any commutative ring) gives
\[
(XI-A)\sum_{k=0}^{n-1}B_kX^k=\chi_A(X)\,I=\sum_{k=0}^nc_kX^kI.
\]
Comparing the coefficients of $X^k$: $B_{k-1}-AB_k=c_kI$ for $0\le k\le n$, with $B_{-1}=B_n=0$. Multiply the equation
for $X^k$ on the left by $A^k$ and add up over $k$:
\[
\sum_{k=0}^n\bigl(A^kB_{k-1}-A^{k+1}B_k\bigr)=\sum_{k=0}^nc_kA^k=\chi_A(A).
\]
The left side telescopes to $A^0B_{-1}-A^{n+1}B_n=0$.
\end{proof}

The tempting one-line ``proof'', substituting $X=A$ in $\det(XI-A)$ to get $\det(A-A)=0$, is wrong: it confuses the
matrix $\chi_A(A)$ with the scalar $\det(0)$. The proof above is the correct version of that idea: the substitution is done
inside the identity $(XI-A)B(X)=\chi_A(X)I$, where it is legitimate because the powers of $A$ are placed on the left.

The polynomials $P$ with $P(A)=0$ form an ideal of $\mathbb K[X]$, which is non-zero by Cayley--Hamilton. Its monic
generator is the \emph{minimal polynomial} $\mu_A$: the monic polynomial of least degree killing $A$, which divides every
polynomial killing $A$ (Euclidean division: the remainder also kills $A$ and has smaller degree, so it is $0$). By
Cayley--Hamilton, $\mu_A$ divides $\chi_A$. Conversely every eigenvalue is a root of $\mu_A$: if $Av=\lambda v$ with $v\ne0$ then
$0=\mu_A(A)v=\mu_A(\lambda)v$. So $\mu_A$ and $\chi_A$ have the same roots, possibly with different multiplicities.

\section{Generalised eigenspaces}

Diagonalisation fails when an eigenvalue has fewer independent eigenvectors than its multiplicity, as for
$\begin{pmatrix}0&1\\0&0\end{pmatrix}$. The remedy is to enlarge eigenspaces to \emph{generalised} eigenspaces. The engine is
a lemma about coprime polynomials.

\begin{lemma}[kernel lemma]\label{linear:lem:kernels}
If $P=P_1\cdots P_k$ with $P_1,\dots,P_k$ pairwise coprime polynomials, then
$\ker P(A)=\ker P_1(A)\oplus\dots\oplus\ker P_k(A)$, and the projections onto the summands are polynomials in $A$.
\end{lemma}
\begin{proof}
Let $Q_i=P/P_i$. The $Q_i$ have no common factor (a prime factor of all of them would divide some $P_j$ and also $Q_j$,
which is a product of polynomials coprime to $P_j$), so by B\'ezout's theorem in $\mathbb K[X]$ there are $U_i$ with
$\sum_iU_iQ_i=1$. Let $E_i=U_i(A)Q_i(A)$; then $\sum_iE_i=I$. For $v\in\ker P(A)$ write $v=\sum_iE_iv$; since
$P_i(A)E_iv=U_i(A)P(A)v=0$, $E_iv\in\ker P_i(A)$. For uniqueness, suppose $\sum_iw_i=0$ with $w_i\in\ker P_i(A)$. For $j\ne i$,
$Q_j$ is divisible by $P_i$, so $E_jw_i=0$; hence $w_i=\sum_jE_jw_i=E_iw_i=E_i\bigl(-\sum_{l\ne i}w_l\bigr)=0$, since $E_i$ kills
each $w_l$, $l\ne i$, by the same remark. So the sum is direct and $E_i$ restricted to $\ker P(A)$ is the projection on
the $i$-th summand.
\end{proof}

\begin{theorem}[generalised eigenspaces]\label{linear:thm:geneigen}
Suppose $\chi_A=\prod_{i=1}^k(X-\lambda_i)^{m_i}$ with distinct $\lambda_i$. Let $V_i=\ker(A-\lambda_i)^{m_i}$, the
\emph{generalised eigenspace} of $\lambda_i$. Then $V=V_1\oplus\dots\oplus V_k$, each $V_i$ is invariant, $\dim V_i=m_i$, and
$A|_{V_i}-\lambda_i$ is nilpotent. Moreover $V_i=\ker(A-\lambda_i)^{n'}$ for every $n'\ge m_i$.
\end{theorem}
\begin{proof}
By Cayley--Hamilton $\ker\chi_A(A)=V$, and the factors $(X-\lambda_i)^{m_i}$ are pairwise coprime, so Lemma~\ref{linear:lem:kernels}
gives the decomposition. Each $V_i$ is the kernel of a polynomial in $A$, so it is invariant, and $(A-\lambda_i)^{m_i}=0$ on
$V_i$. The operator $A|_{V_i}$ has the single eigenvalue $\lambda_i$ (an eigenvector $v\in V_i$ for $\mu$ satisfies
$0=(A-\lambda_i)^{m_i}v=(\mu-\lambda_i)^{m_i}v$), so $\chi_{A|_{V_i}}=(X-\lambda_i)^{d_i}$ with $d_i=\dim V_i$. In a basis adapted to the
decomposition $A$ is block diagonal, so $\chi_A=\prod_i(X-\lambda_i)^{d_i}$, and $d_i=m_i$ by unique factorisation. Finally
$\ker(A-\lambda_i)^{n'}$ is invariant and $A$ has only the eigenvalue $\lambda_i$ on it, so it meets $V_j$, $j\ne i$, only in
$0$ ($A-\lambda_i$ is invertible on $V_j$); by the decomposition it is contained in $V_i$, and it contains $V_i$ for $n'\ge
m_i$.
\end{proof}

\begin{corollary}[diagonalisability]\label{linear:cor:diag}
$A$ is diagonalisable over $\mathbb K$ if and only if $\mu_A$ splits with simple roots, if and only if $A$ is killed by some
polynomial that splits with simple roots.
\end{corollary}
\begin{proof}
If $A$ is diagonalisable with distinct eigenvalues $\lambda_1,\dots,\lambda_k$, then $\prod_i(X-\lambda_i)$ kills $A$ (check on
a basis of eigenvectors). Conversely, if $P=\prod_i(X-\lambda_i)$ with distinct $\lambda_i$ kills $A$, the kernel lemma gives
$V=\ker P(A)=\bigoplus_i\ker(A-\lambda_i)$, a sum of eigenspaces. The middle condition is equivalent because $\mu_A$ divides
every polynomial killing $A$.
\end{proof}

\begin{theorem}[Dunford--Jordan--Chevalley decomposition]\label{linear:thm:dunford}
If $\chi_A$ splits, there is a unique pair $(D,N)$ of operators with $A=D+N$, $D$ diagonalisable, $N$ nilpotent and $DN=ND$.
Both $D$ and $N$ are polynomials in $A$.
\end{theorem}
\begin{proof}
With the notation of Theorem~\ref{linear:thm:geneigen}, let $D$ be $\lambda_i$ on $V_i$, that is, $D=\sum_i\lambda_iE_i$ with the
projections $E_i$ of Lemma~\ref{linear:lem:kernels}, which are polynomials in $A$. Then $D$ is a polynomial in $A$, it is
diagonalisable, $N=A-D$ is a polynomial in $A$, so $N$ and $D$ commute, and $N$ is $A-\lambda_i$ on $V_i$, which is
nilpotent. For uniqueness, let $A=D'+N'$ be another such decomposition. Then $D'$ commutes with $A=D'+N'$, hence with
every polynomial in $A$, in particular with $D$ and $N$. Two commuting diagonalisable operators are simultaneously
diagonalisable (Theorem~\ref{linear:thm:commuting} below), so $D-D'$ is diagonalisable; and $N'-N$ is nilpotent, as a difference
of commuting nilpotent operators (expand $(N'-N)^{2n}$ with the binomial formula: each term contains $N^n$ or $N'^n$).
Since $D-D'=N'-N$ is both diagonalisable and nilpotent, its eigenvalues are all $0$, so it is $0$.
\end{proof}

\section{Nilpotent operators and Young diagrams}\label{linear:sec:nilpotent}

By Theorem~\ref{linear:thm:geneigen}, the whole structure of an operator with split characteristic polynomial comes down to its
eigenvalues and the nilpotent operators $A-\lambda_i$ on the generalised eigenspaces. So let $N$ be nilpotent on $V$, of
dimension $n$, and let $s$ be the least integer with $N^s=0$. The kernels of its powers increase,
\[
\{0\}=\ker N^0\subsetneq\ker N\subsetneq\ker N^2\subsetneq\dots\subsetneq\ker N^s=V,
\]
strictly: if $\ker N^k=\ker N^{k+1}$ for some $k<s$, then $N^{k+2}v=0$ gives $Nv\in\ker N^{k+1}=\ker N^k$, so $v\in\ker
N^{k+1}$, and inductively the kernels would stop growing before reaching $V$. Let
\[
r_k=\dim\ker N^k-\dim\ker N^{k-1}\qquad(k\ge1).
\]

\begin{lemma}\label{linear:lem:decreasing}
$N$ induces injective linear maps
\[
\ker N^{k+1}/\ker N^k\longrightarrow\ker N^k/\ker N^{k-1}.
\]
Consequently $r_1\ge r_2\ge\dots\ge r_s>0$
and $r_1+\dots+r_s=n$.
\end{lemma}
\begin{proof}
$N$ maps $\ker N^{k+1}$ into $\ker N^k$ and $\ker N^k$ into $\ker N^{k-1}$, so it induces a map on the quotients. If $Nv\in\ker
N^{k-1}$ for $v\in\ker N^{k+1}$, then $N^kv=0$, so $v\in\ker N^k$: the induced map is injective, and $r_{k+1}\le r_k$. The sum
telescopes to $\dim\ker N^s=n$.
\end{proof}

So $(r_1,\dots,r_s)$ is a partition of $n$. Draw its \emph{Young diagram} with columns of heights $r_1,r_2,\dots$ from left
to right (Figure~\ref{linear:fig:young}); its rows have lengths $m_1\ge m_2\ge\dots$, the \emph{conjugate} partition. The Jordan
block of size $m$ is the $m\times m$ matrix $J_m$ with ones just above the diagonal and zeros elsewhere: $J_me_1=0$ and
$J_me_j=e_{j-1}$.

\begin{figure}[ht]
\centering
\fig{linear-ch03-young}
\caption{The Young diagram of a nilpotent operator with $\dim\ker N^k=4,7,8,9$ for $k=1,2,3,4$. Each row is a Jordan
chain $u,Nu,N^2u,\dots$ read from right to left, ending in the kernel of $N$ (first column); here the Jordan blocks have
sizes $4,2,2,1$.}
\label{linear:fig:young}
\end{figure}

\begin{theorem}[Jordan form of a nilpotent operator]\label{linear:thm:jordannil}
There is a basis of $V$ in which $N$ is block diagonal, $N=\operatorname{diag}(J_{m_1},\dots,J_{m_t})$, where $m_1\ge\dots\ge
m_t$ are the row lengths of the Young diagram of $(r_k)$. The sizes $m_i$ are uniquely determined by $N$: the number of
blocks of size at least $k$ is $r_k$.
\end{theorem}
\begin{proof}
\emph{Existence.} Write $K_k=\ker N^k$. We choose vectors level by level, from $k=s$ down to $k=1$: at level $k$ we choose
$u_{k,1},\dots,u_{k,t_k}\in K_k$ such that the classes in $K_k/K_{k-1}$ of
\[
N^{j-k}u_{j,i}\quad(j>k,\ 1\le i\le t_j)\qquad\text{and}\qquad u_{k,1},\dots,u_{k,t_k}
\]
form a basis of $K_k/K_{k-1}$. This is possible: the classes of the $N^{j-k}u_{j,i}$, $j>k$, are linearly independent in
$K_k/K_{k-1}$, because they are the images under the injective map of Lemma~\ref{linear:lem:decreasing} of the classes of
$N^{j-k-1}u_{j,i}$, which are part of a basis of $K_{k+1}/K_k$ by the previous step; so we can complete them into a basis
(Theorem~\ref{linear:thm:dim}). Now consider all the vectors $N^lu_{k,i}$ with $0\le l<k$. For each level $k'$, the ones that lie in
$K_{k'}\setminus K_{k'-1}$, namely those with $k-l=k'$, have classes forming a basis of $K_{k'}/K_{k'-1}$ by construction.
A family of vectors in $V$ whose members, sorted by level, give bases of all the quotients $K_{k'}/K_{k'-1}$ is a basis of $V$:
there are $\sum_{k'}r_{k'}=n$ of them, and they are independent, since in a vanishing combination the terms of the highest
level present would give a vanishing combination of classes in that quotient. For each $(k,i)$ the chain
$N^{k-1}u_{k,i},\dots,Nu_{k,i},u_{k,i}$, in this order, spans an invariant subspace on which $N$ acts by $J_k$. So $N$ is block
diagonal with $t_k$ blocks of size $k$ for each $k$.

\emph{Counting.} In the block form, $\dim\ker N^k=\sum_i\min(m_i,k)$, since $\dim\ker J_m^k=\min(m,k)$. Hence $r_k$ is the
number of $i$ with $m_i\ge k$: the column heights of the diagram are $r_k$ and the rows are the $m_i$. In particular the
block sizes are determined by the numbers $\dim\ker N^k$, which do not depend on the basis.
\end{proof}

Combining the generalised eigenspaces with this theorem gives the Jordan normal form.

\begin{corollary}[Jordan normal form]\label{linear:cor:jordan}
If $\chi_A$ splits, $A$ is similar to a block diagonal matrix of Jordan blocks $\lambda I_m+J_m$, unique up to the order of
the blocks. The number of blocks of size at least $k$ for the eigenvalue $\lambda$ is
$\operatorname{rank}(A-\lambda)^{k-1}-\operatorname{rank}(A-\lambda)^k$. Two such matrices are similar if and only if
they have the same eigenvalues and $\operatorname{rank}(A-\lambda)^k=\operatorname{rank}(B-\lambda)^k$ for all $\lambda$ and $k$.
\end{corollary}
\begin{proof}
Apply Theorem~\ref{linear:thm:jordannil} to $A-\lambda_i$ on each $V_i$. On $V_j$, $j\ne i$, the operator $A-\lambda_i$ is invertible, so
$\dim\ker(A-\lambda_i)^k=\dim\ker(A-\lambda_i)^k|_{V_i}$, and the rank formula is rank--nullity. Similar matrices have the same
ranks of all $(A-\lambda)^k$; conversely, the ranks determine the Jordan form, and two matrices with the same Jordan form
are similar to it, hence to each other.
\end{proof}

\begin{corollary}\label{linear:cor:transpose}
Every complex square matrix is similar to its transpose.
\end{corollary}
\begin{proof}
$\operatorname{rank}(A^{\mathsf T}-\lambda)^k=\operatorname{rank}((A-\lambda)^k)^{\mathsf T}=\operatorname{rank}(A-\lambda)^k$
(Theorem~\ref{linear:thm:transpose}), and $A$, $A^{\mathsf T}$ have the same characteristic polynomial.
\end{proof}

\begin{remark}
Over a field where $\chi_A$ does not split (such as $\R$ or $\Q$), the theory of modules over the principal ideal
domain $\mathbb K[X]$ gives the \emph{rational canonical form}, and the Jordan form appears as a special case; this is
the route taken in Tome~2, Chapter~7. Corollary~\ref{linear:cor:transpose} holds over every field, by that theory. The proof here
has the advantage of producing the Jordan basis by a linear-algebra algorithm, and of explaining the shape of the Young
diagram: the columns are measured by kernels, the rows are chains.
\end{remark}

The Jordan form is beautiful but numerically treacherous. The matrix $J_n$ has the single eigenvalue $0$, but changing its
bottom-left entry from $0$ to $\varepsilon$ gives a matrix with characteristic polynomial $X^n-\varepsilon$, whose eigenvalues
are the $n$-th roots of $\varepsilon$, all of modulus $\varepsilon^{1/n}$ (Figure~\ref{linear:fig:perturb}). For $n=12$ and
$\varepsilon=10^{-6}$, a perturbation at the level of rounding errors in single precision moves the eigenvalues by more than $0.3$.
The Jordan form is not a continuous function of the matrix, and computers never compute it in floating point; the
singular value decomposition of Appendix~F and Schur's triangular form, which are stable, take its place.

\begin{figure}[ht]
\centering
\fig{linear-ch03-perturb}
\caption{The eigenvalues of $J_{12}+10^{-6}E_{12,1}$: a perturbation of size $10^{-6}$ moves the twelve-fold eigenvalue $0$
to a circle of radius $10^{-1/2}\approx0.32$.}
\label{linear:fig:perturb}
\end{figure}

\section{Commuting operators}

\begin{theorem}\label{linear:thm:commuting}
Let $\mathcal F$ be a family of pairwise commuting operators on $V$.
\begin{enumerate}[label=(\alph*)]
\item If every member of $\mathcal F$ is diagonalisable, they are simultaneously diagonalisable.
\item If every characteristic polynomial splits (for instance over $\C$), the members of $\mathcal F$ have a common
eigenvector and are simultaneously triangularisable.
\end{enumerate}
\end{theorem}
\begin{proof}
(a) By induction on $\dim V$. If every member of $\mathcal F$ is a scalar, any basis works. Otherwise some $A\in\mathcal F$
has at least two distinct eigenvalues, and $V=\bigoplus_\lambda\ker(A-\lambda)$ with each summand of smaller dimension. Every
$B\in\mathcal F$ preserves each eigenspace of $A$: if $Av=\lambda v$ then $A(Bv)=BAv=\lambda Bv$. The restriction of $B$ to an
invariant subspace is still diagonalisable (it is killed by $\mu_B$, which has simple roots; use Corollary~\ref{linear:cor:diag}).
By induction the restrictions of $\mathcal F$ to each eigenspace are simultaneously diagonalisable, and the union of the
bases works.

(b) Common eigenvector: choose a non-zero invariant subspace $W$ of minimal dimension among the subspaces invariant under
every member of $\mathcal F$. If some $A\in\mathcal F$ is not scalar on $W$, pick an eigenvalue $\lambda$ of $A|_W$ (its
characteristic polynomial divides $\chi_A$, so it splits); then $\ker(A-\lambda)\cap W$ is non-zero, smaller than $W$, and
invariant under $\mathcal F$ by the computation in (a): a contradiction. So every member is scalar on $W$, and any non-zero
$w\in W$ is a common eigenvector. Triangularisation follows by induction exactly as in Theorem~\ref{linear:thm:schur}, using
the common eigenvector and the induced commuting operators on $V/\mathbb Kw$, whose characteristic polynomials still split
by~\eqref{linear:eq:charfactor}.
\end{proof}

\section{Functions of a matrix}

Polynomials in $A$ are easy: if $A=PJP^{-1}$, then $f(A)=Pf(J)P^{-1}$, and $f(J)$ is block diagonal. On one Jordan block
$\lambda I+J_m$, Taylor's formula for the polynomial $f$, $f(\lambda+x)=\sum_jf^{(j)}(\lambda)x^j/j!$, with $x=J_m$ and $J_m^m=0$,
gives
\begin{equation}\label{linear:eq:fjordan}
\begin{aligned}
f(\lambda I+J_m)&=\sum_{j=0}^{m-1}\frac{f^{(j)}(\lambda)}{j!}J_m^j\\
&=\begin{pmatrix}f(\lambda)&f'(\lambda)&\frac{f''(\lambda)}2&\cdots\\&f(\lambda)&f'(\lambda)&\ddots\\&&\ddots&\ddots\\&&&f(\lambda)\end{pmatrix}.
\end{aligned}
\end{equation}
So $f(A)$ only depends on the values of $f$ and its derivatives up to order $m_\lambda-1$ at each eigenvalue, $m_\lambda$ the size
of the largest block for $\lambda$: exactly the multiplicity of $\lambda$ as a root of $\mu_A$. This lets us define $f(A)$ for
any function $f$ that is smooth near the eigenvalues, and agrees with power series when they converge. The most
important case is the exponential.

\begin{proposition}[the matrix exponential]\label{linear:prop:exp}
For $A\in M_n(\C)$ the series $e^A=\sum_{k\ge0}A^k/k!$ converges absolutely (entrywise). It satisfies $e^{A+B}=e^Ae^B$ when
$AB=BA$, $e^{PAP^{-1}}=Pe^AP^{-1}$, $\det e^A=e^{\operatorname{tr}A}$, and $\frac{\mathrm d}{\mathrm dt}e^{tA}=Ae^{tA}$. In
particular $e^A$ is always invertible, with inverse $e^{-A}$.
\end{proposition}
\begin{proof}
Let $\|A\|=n\max|a_{ij}|$. Then $\|AB\|\le\|A\|\|B\|$ (each entry of $AB$ is a sum of $n$ products, bounded by
$n\max|a|\max|b|$), so the entries of $A^k/k!$ are bounded by $\|A\|^k/k!$, and the series converges absolutely. For
commuting $A,B$, the binomial formula $(A+B)^m=\sum_k\binom mkA^kB^{m-k}$ holds, and the Cauchy product of the two absolutely
convergent series gives $e^Ae^B=\sum_m\sum_{k}\frac{A^kB^{m-k}}{k!(m-k)!}=e^{A+B}$. Conjugation passes through each partial
sum. For the determinant, triangularise $A=PTP^{-1}$: $e^T$ is triangular with diagonal $e^{t_{ii}}$ (each $T^k$ is triangular
with diagonal $t_{ii}^k$), so $\det e^A=\det e^T=\prod e^{t_{ii}}=e^{\operatorname{tr}A}$. The derivative: $e^{(t+h)A}=e^{tA}e^{hA}$
and $(e^{hA}-I)/h=A+hA^2\sum_{k\ge2}h^{k-2}A^{k-2}/k!\to A$.
\end{proof}

For example $\exp\begin{pmatrix}0&-t\\t&0\end{pmatrix}=\begin{pmatrix}\cos t&-\sin t\\\sin t&\cos t\end{pmatrix}$
(Exercise~\ref{linear:ex:rotation}): the exponential turns an infinitesimal rotation into a rotation. Appendix~L uses the
exponential to solve linear differential equations, and Tome~8 builds Lie theory on it.

\begin{beauty}
A single number, the dimension of a kernel, seems a crude invariant. Yet the sequence $\dim\ker N,\dim\ker N^2,\dots$
determines a nilpotent operator completely up to change of basis, and its successive differences arrange themselves into a
Young diagram whose columns are measured by kernels and whose rows are Jordan chains. The classification of nilpotent
matrices is the classification of partitions; the same diagrams index the irreducible representations of the symmetric
group (Tome~9), and this is no coincidence but the beginning of a vast dictionary.
\end{beauty}

\section{Exercises}

\begin{exercise}[$\star$]\label{linear:ex:abba}
Show that $AB$ and $BA$ have the same characteristic polynomial for $A,B\in M_n(\mathbb K)$. (First suppose $A$ invertible.)
For $A\in M_{m,n}$ and $B\in M_{n,m}$ with $m\le n$, show that $\chi_{BA}=X^{n-m}\chi_{AB}$.
\end{exercise}

\begin{exercise}[$\star$]\label{linear:ex:companion}
Show that the companion matrix of $P=X^n+a_{n-1}X^{n-1}+\dots+a_0$ (ones below the diagonal, last column
$-a_0,\dots,-a_{n-1}$) has characteristic polynomial and minimal polynomial both equal to $P$.
\end{exercise}

\begin{exercise}[$\star$]\label{linear:ex:nilcount}
How many similarity classes of $n\times n$ complex matrices $N$ satisfy $N^2=0$? $N^3=0$ and $n=6$?
\end{exercise}

\begin{exercise}[$\star\star$]\label{linear:ex:jordanfind}
Find the Jordan form of a $7\times7$ matrix $A$ with $\chi_A=(X-2)^5(X+1)^2$, $\operatorname{rank}(A-2I)=4$,
$\operatorname{rank}(A-2I)^2=3$, $\operatorname{rank}(A-2I)^3=2$ and $\operatorname{rank}(A+I)=6$. What is $\mu_A$?
\end{exercise}

\begin{exercise}[$\star\star$]\label{linear:ex:finiteorder}
Show that if $A\in M_n(\C)$ satisfies $A^k=I$ for some $k\ge1$, then $A$ is diagonalisable. Is this true over
$\mathbb F_2$?
\end{exercise}

\begin{exercise}[$\star\star$]\label{linear:ex:nilptrace}
Over a field of characteristic $0$, show that $A$ is nilpotent if and only if $\operatorname{tr}(A^k)=0$ for $k=1,\dots,n$.
\end{exercise}

\begin{exercise}[$\star\star$]\label{linear:ex:sqrt}
Show that $J_n$ ($n\ge2$) has no square root, while $\lambda I+J_n$ has one for every $\lambda\ne0$ in $\C$.
\end{exercise}

\begin{exercise}[$\star\star$]\label{linear:ex:commutant}
Let $A\in M_n(\C)$. Show that the matrices commuting with $A$ form a subspace of dimension at least $n$, with equality
when $A$ has $n$ distinct eigenvalues, and that in that case they are exactly the polynomials in $A$.
\end{exercise}

\begin{exercise}[$\star$]\label{linear:ex:rotation}
Compute $\exp\begin{pmatrix}0&-t\\t&0\end{pmatrix}$ and $\exp(t(\lambda I+J_3))$.
\end{exercise}

\begin{exercise}[$\star\star$]\label{linear:ex:minpolyjordan}
Show that the minimal polynomial of $A$ (with split $\chi_A$) is $\prod_\lambda(X-\lambda)^{s_\lambda}$, where $s_\lambda$ is the size of
the largest Jordan block for $\lambda$. Deduce that $\mu_A=\chi_A$ if and only if each eigenvalue has a single Jordan block.
\end{exercise}

\begin{exercise}[$\star\star$]\label{linear:ex:expsum}
Find $2\times2$ matrices $A,B$ with $e^{A+B}\ne e^Ae^B$. Find $A\ne0$ with $e^A=I$.
\end{exercise}

\begin{exercise}[$\star\star\star$]\label{linear:ex:symprod}
(a) Let $R$ be the $m\times m$ matrix with ones on the anti-diagonal. Show that $RJ_mR=J_m^{\mathsf T}$. (b) Deduce that for
every $A\in M_n(\C)$ there is an invertible \emph{symmetric} $S$ with $A^{\mathsf T}=S^{-1}AS$. (c) Deduce that every complex
square matrix is a product of two symmetric matrices, one of them invertible (Frobenius, 1910).
\end{exercise}

\section{Hints}

\begin{hint}{linear:ex:abba}
$AB=A(BA)A^{-1}$. For the general case consider $A+tI$, or the block matrices $\begin{pmatrix}AB&0\\B&0\end{pmatrix}$ and
$\begin{pmatrix}0&0\\B&BA\end{pmatrix}$, conjugate by $\begin{pmatrix}I&A\\0&I\end{pmatrix}$.
\end{hint}
\begin{hint}{linear:ex:companion}
Expand $\det(XI-C)$ along the last column; for the minimal polynomial, look at $e_1,Ce_1,\dots,C^{n-1}e_1$.
\end{hint}
\begin{hint}{linear:ex:nilcount}
Partitions of $n$ with parts at most $2$, resp.\ partitions of $6$ with parts at most $3$.
\end{hint}
\begin{hint}{linear:ex:jordanfind}
Compute the $r_k=\operatorname{rank}(A-2)^{k-1}-\operatorname{rank}(A-2)^k$ and draw the Young diagram.
\end{hint}
\begin{hint}{linear:ex:finiteorder}
$X^k-1$ has simple roots in $\C$. Over $\mathbb F_2$, try $\begin{pmatrix}1&1\\0&1\end{pmatrix}$.
\end{hint}
\begin{hint}{linear:ex:nilptrace}
Write the traces with the eigenvalues; a Vandermonde system shows that the distinct non-zero eigenvalues have zero
multiplicity.
\end{hint}
\begin{hint}{linear:ex:sqrt}
If $B^2=J_n$, then $B$ is nilpotent; compare $\dim\ker B$ and $\dim\ker B^2$. For $\lambda\ne0$, use \eqref{linear:eq:fjordan} with a
branch of $\sqrt{\phantom{x}}$ near $\lambda$, or the binomial series.
\end{hint}
\begin{hint}{linear:ex:commutant}
For diagonal $A$ with distinct entries the commutant is the diagonal matrices. In general use the Jordan form and count
the matrices commuting with each block.
\end{hint}
\begin{hint}{linear:ex:rotation}
Diagonalise over $\C$, or sum the series, separating even and odd powers.
\end{hint}
\begin{hint}{linear:ex:minpolyjordan}
$(J_m)^{k}=0$ if and only if $k\ge m$.
\end{hint}
\begin{hint}{linear:ex:expsum}
Try $A=E_{12}$, $B=E_{21}$. For the second question, think of rotations.
\end{hint}
\begin{hint}{linear:ex:symprod}
(b) Use the Jordan form $A=PJP^{-1}$ and a block diagonal $R$. (c) $AS=SA^{\mathsf T}$, so $(AS)^{\mathsf T}=AS$.
\end{hint}

\section{Solutions}

\begin{solution}{linear:ex:abba}
\idea Similar matrices have the same characteristic polynomial, and $AB$ is similar to $BA$ when $A$ is invertible; the
general case follows by continuity or by an exact block trick.
\writeup If $A$ is invertible, $AB=A(BA)A^{-1}$. In general, with $m,n$ arbitrary, let $M=\begin{pmatrix}AB&0\\B&0\end{pmatrix}$ and
$M'=\begin{pmatrix}0&0\\B&BA\end{pmatrix}$, of size $m+n$, and $P=\begin{pmatrix}I_m&A\\0&I_n\end{pmatrix}$, invertible. A direct computation gives
$MP=\begin{pmatrix}AB&ABA\\B&BA\end{pmatrix}=PM'$, so $M$ and $M'$ are similar. Their characteristic polynomials are
$\chi_{AB}X^n$ and $X^m\chi_{BA}$ (block triangular). So $X^n\chi_{AB}=X^m\chi_{BA}$; for $m\le n$, $\chi_{BA}=X^{n-m}\chi_{AB}$, and for
$m=n$ they are equal. This argument works over any commutative ring, with no division.
\end{solution}

\begin{solution}{linear:ex:companion}
\idea The companion matrix is built so that $Ce_j=e_{j+1}$: the vector $e_1$ generates everything, which forces the
minimal polynomial to have full degree.
\writeup Expanding $\det(XI-C)$ along the last column (or by induction on $n$, expanding along the first row) gives
$X^n+a_{n-1}X^{n-1}+\dots+a_0=P$. For the minimal polynomial: $Ce_j=e_{j+1}$ for $j<n$, so $e_1,Ce_1,\dots,C^{n-1}e_1$ is the basis
$e_1,\dots,e_n$. If $Q$ of degree $d<n$ killed $C$, then $Q(C)e_1=0$ would be a non-trivial linear relation between
$e_1,\dots,C^de_1$. So $\deg\mu_C=n$, and as $\mu_C$ divides $\chi_C=P$, both are monic of degree $n$: $\mu_C=P$.
\end{solution}

\begin{solution}{linear:ex:nilcount}
\idea Similarity classes of nilpotent matrices are partitions of $n$ (the block sizes); $N^k=0$ means all blocks have size
at most $k$.
\writeup $N^2=0$: partitions of $n$ into parts $1$ and $2$, determined by the number $j$ of parts equal to $2$,
$0\le j\le\lfloor n/2\rfloor$: there are $\lfloor n/2\rfloor+1$ classes (they are distinguished by the rank $j$). $N^3=0$, $n=6$:
partitions of $6$ into parts at most $3$: $3{+}3$, $3{+}2{+}1$, $3{+}1{+}1{+}1$, $2{+}2{+}2$, $2{+}2{+}1{+}1$,
$2{+}1{+}1{+}1{+}1$, $1{+}1{+}1{+}1{+}1{+}1$: seven classes.
\end{solution}

\begin{solution}{linear:ex:jordanfind}
\idea The number of blocks of size $\ge k$ is $r_k=\operatorname{rank}(A-\lambda)^{k-1}-\operatorname{rank}(A-\lambda)^k$; on the
generalised eigenspace of $2$ the ranks must eventually drop by $5$ from $7$.
\writeup For $\lambda=2$: ranks $7,4,3,2$ for $k=0,1,2,3$, and $\operatorname{rank}(A-2)^k=2$ for $k\ge3$ (it cannot drop
below the dimension $2$ of the other generalised eigenspace, on which $A-2$ is invertible, and $3+1+1=5$ already accounts
for the multiplicity). So $r_1=3$, $r_2=1$, $r_3=1$: three blocks, one of size $\ge2$, one of size $\ge3$; sizes $3,1,1$
(total $5$). For $\lambda=-1$: $\operatorname{rank}(A+I)=6$ means one block, of size $2$. So
$A\sim\operatorname{diag}(2I_3+J_3,\ 2,\ 2,\ -I_2+J_2)$, and $\mu_A=(X-2)^3(X+1)^2$ (Exercise~\ref{linear:ex:minpolyjordan}).
\end{solution}

\begin{solution}{linear:ex:finiteorder}
\idea A polynomial with simple roots killing $A$ forces diagonalisability.
\writeup $X^k-1$ kills $A$, and it has $k$ distinct roots in $\C$, the $k$-th roots of unity (its derivative $kX^{k-1}$ has
no common root with it). By Corollary~\ref{linear:cor:diag}, $A$ is diagonalisable. Over $\mathbb F_2$,
$A=\begin{pmatrix}1&1\\0&1\end{pmatrix}$ satisfies $A^2=I$ but is not diagonalisable: its only eigenvalue is $1$ and $A\ne I$. Here
$X^2-1=(X-1)^2$ in characteristic $2$, so the root is not simple.
\end{solution}

\begin{solution}{linear:ex:nilptrace}
\idea Traces of powers are power sums of the eigenvalues; power sums determine which non-zero values occur, by a
Vandermonde determinant.
\writeup Over an extension where $\chi_A$ splits (or over $\C$ if $\mathbb K\subseteq\C$), let the distinct non-zero eigenvalues
be $\mu_1,\dots,\mu_r$ with multiplicities $n_1,\dots,n_r\ge1$. Then $\operatorname{tr}(A^k)=\sum_jn_j\mu_j^k$. If these vanish for
$k=1,\dots,r$ ($r\le n$), the vector $(n_1\mu_1,\dots,n_r\mu_r)$ is killed by the matrix $(\mu_j^{k-1})_{k,j}$, a Vandermonde matrix
with distinct nodes, hence invertible. So $n_j\mu_j=0$; in characteristic $0$, $n_j\ne0$ as an element of the field, so
$\mu_j=0$, a contradiction unless $r=0$. So all eigenvalues are $0$, $\chi_A=X^n$, and $A^n=0$ by Cayley--Hamilton. The
converse is clear: a nilpotent matrix has nilpotent powers, all with eigenvalues $0$ and trace $0$. (In characteristic
$p$ the statement fails: $I_p$ has $\operatorname{tr}(I_p^k)=p=0$.)
\end{solution}

\begin{solution}{linear:ex:sqrt}
\idea A square root of a nilpotent is nilpotent, and kernels of powers of a nilpotent grow by at most the size of the
first kernel; for $J_n$ they grow by $1$ at each step, too slowly.
\writeup If $B^2=J_n$ then $B^{2n}=0$, so $B$ is nilpotent, and $B\ne0$. Then $\dim\ker B\ge1$ and, by
Lemma~\ref{linear:lem:decreasing}, $\dim\ker B^2\le2\dim\ker B$, while $\dim\ker B^2=\dim\ker J_n=1$ forces $\dim\ker B=1$. Then
$B$ has a single Jordan block of size $n$, so $\dim\ker B^2=\min(2,n)=2$ for $n\ge2$: a contradiction. For $\lambda\ne0$,
choose a square root $\sigma$ of $\lambda$ and set $B=\sigma\sum_{j=0}^{n-1}\binom{1/2}{j}(J_n/\lambda)^j$, the binomial series of
$\sigma\sqrt{1+x}$ at $x=J_n/\lambda$, truncated because $J_n^n=0$. The identity $\bigl(\sum_j\binom{1/2}jx^j\bigr)^2=1+x$ holds for
power series, hence modulo $x^n$, hence for $x=J_n/\lambda$; so $B^2=\lambda(I+J_n/\lambda)=\lambda I+J_n$.
\end{solution}

\begin{solution}{linear:ex:commutant}
\idea Commuting is a linear condition, invariant under conjugation, so we may assume $A$ in Jordan form; then count.
\writeup The commutant $\mathcal C(A)=\{X:AX=XA\}$ is a subspace, and $\mathcal C(PAP^{-1})=P\,\mathcal C(A)P^{-1}$, so we may take $A=J$
in Jordan form. The polynomials in $J_m$ ($I,J_m,\dots,J_m^{m-1}$, independent) commute with $J_m$; placing such a polynomial
in each diagonal block gives $\sum_im_i=n$ independent matrices commuting with $J$. So $\dim\mathcal C(A)\ge n$. If $A$ is
diagonal with distinct entries $\lambda_i$, $AX=XA$ reads $(\lambda_i-\lambda_j)x_{ij}=0$, so $X$ is diagonal: $\dim=n$. The
polynomials in $A$ are diagonal, and by Lagrange interpolation every diagonal matrix $\operatorname{diag}(d_i)$ equals
$P(A)$ for the polynomial with $P(\lambda_i)=d_i$. So the commutant is $\mathbb K[A]$.
\end{solution}

\begin{solution}{linear:ex:rotation}
\idea Split the exponential series into even and odd powers, as for Euler's formula.
\writeup With $K=\begin{pmatrix}0&-1\\1&0\end{pmatrix}$, $K^2=-I$, so $(tK)^{2j}=(-1)^jt^{2j}I$ and $(tK)^{2j+1}=(-1)^jt^{2j+1}K$.
Summing, $e^{tK}=\cos t\,I+\sin t\,K=\begin{pmatrix}\cos t&-\sin t\\\sin t&\cos t\end{pmatrix}$. For $\lambda I+J_3$: $\lambda I$ and
$J_3$ commute, so $e^{t(\lambda I+J_3)}=e^{\lambda t}e^{tJ_3}=e^{\lambda t}\bigl(I+tJ_3+\tfrac{t^2}2J_3^2\bigr)$, the matrix with
$e^{\lambda t}$ on the diagonal, $te^{\lambda t}$ above it and $\frac{t^2}2e^{\lambda t}$ in the corner. This agrees with
\eqref{linear:eq:fjordan} for $f(x)=e^{tx}$.
\end{solution}

\begin{solution}{linear:ex:minpolyjordan}
\idea A polynomial kills a block diagonal matrix if and only if it kills each block, and $(\lambda-\mu)$-shifts are
invertible on blocks of other eigenvalues.
\writeup $P(A)=0$ if and only if $P(\lambda I+J_m)=0$ for every block. By~\eqref{linear:eq:fjordan}, $P(\lambda I+J_m)=0$ if and only if
$P(\lambda)=P'(\lambda)=\dots=P^{(m-1)}(\lambda)=0$, that is, $(X-\lambda)^m$ divides $P$. So the monic polynomial of least degree
killing $A$ is $\prod_\lambda(X-\lambda)^{s_\lambda}$. It equals $\chi_A=\prod_\lambda(X-\lambda)^{m_\lambda}$ if and only if $s_\lambda=m_\lambda$ for every
$\lambda$, that is, the largest block has the full multiplicity, so it is the only block.
\end{solution}

\begin{solution}{linear:ex:expsum}
\idea Without commutation the binomial formula fails; the simplest non-commuting pair is $E_{12},E_{21}$.
\writeup $A=E_{12}$ and $B=E_{21}$ are nilpotent, so $e^A=I+A=\begin{pmatrix}1&1\\0&1\end{pmatrix}$, $e^B=\begin{pmatrix}1&0\\1&1\end{pmatrix}$ and
$e^Ae^B=\begin{pmatrix}2&1\\1&1\end{pmatrix}$. But $A+B=\begin{pmatrix}0&1\\1&0\end{pmatrix}$ has $(A+B)^2=I$, so
$e^{A+B}=\cosh1\,I+\sinh1\,(A+B)$, whose diagonal entries are $\cosh1\approx1.543\ne2$. For the second question,
$A=2\pi K$ with $K$ as in Exercise~\ref{linear:ex:rotation} gives $e^A=\cos2\pi\,I+\sin2\pi\,K=I$.
\end{solution}

\begin{solution}{linear:ex:symprod}
\idea The reversal $R$ conjugates a Jordan block to its transpose and is symmetric; conjugating it by the Jordan basis
keeps it symmetric if we use $P^{-\mathsf T}$ on one side.
\writeup (a) $J_me_j=e_{j-1}$ and $Re_j=e_{m+1-j}$, so $RJ_mRe_j=RJ_me_{m+1-j}=Re_{m-j}=e_{j+1}=J_m^{\mathsf T}e_j$ (with the
convention $e_0=e_{m+1}=0$). (b) Write $A=PJP^{-1}$ with $J$ in Jordan form, and let $R$ be block diagonal with an
anti-diagonal block for each Jordan block; then $R=R^{\mathsf T}=R^{-1}$ and $RJR=J^{\mathsf T}$ (blockwise, using (a) and the fact
that $R$ commutes with $\lambda I$). Then
$A^{\mathsf T}=P^{-\mathsf T}J^{\mathsf T}P^{\mathsf T}=P^{-\mathsf T}RJRP^{\mathsf T}=(P^{-\mathsf T}RP^{-1})A(PRP^{\mathsf T})$. Let
$S=PRP^{\mathsf T}$, which is symmetric and invertible, with $S^{-1}=P^{-\mathsf T}RP^{-1}$. So $A^{\mathsf T}=S^{-1}AS$. (c) From
$SA^{\mathsf T}=AS$ we get $(AS)^{\mathsf T}=S^{\mathsf T}A^{\mathsf T}=SA^{\mathsf T}=AS$: the matrix $T=AS$ is symmetric, and $A=TS^{-1}$ is the
product of two symmetric matrices, $S^{-1}$ being invertible.
\end{solution}

\section*{Sources and further reading}

\begin{itemize}
\item A.~Cayley, \emph{A memoir on the theory of matrices}, Philosophical Transactions of the Royal Society 148 (1858),
17--37.
\item A.-L.~Cauchy, \emph{Sur l'\'equation \`a l'aide de laquelle on d\'etermine les in\'egalit\'es s\'eculaires des
mouvements des plan\`etes}, Exercices de math\'ematiques 4 (1829). Eigenvalues of symmetric matrices are real.
\item T.~Hawkins, \emph{The Mathematics of Frobenius in Context}, Springer, 2013, and \emph{Another look at Cayley and the
theory of matrices}, Archives internationales d'histoire des sciences 27 (1977), 82--112, for the history told above.
\item F.~Brechenmacher, \emph{La controverse de 1874 entre Camille Jordan et Leopold Kronecker}, Revue d'histoire des
math\'ematiques 13 (2007), 187--257.
\item C.~Jordan, \emph{Trait\'e des substitutions et des \'equations alg\'ebriques}, Gauthier-Villars, 1870, \S\S157--159.
\item S.~Axler, \emph{Linear Algebra Done Right}, 4th ed., Springer, 2024, Chapter~8 (generalised eigenspaces and the
Jordan form without determinants).
\item R.~A.~Horn and C.~R.~Johnson, \emph{Matrix Analysis}, 2nd ed., Cambridge University Press, 2013, Chapter~3, including
the symmetric similarity of Exercise~\ref{linear:ex:symprod}.
\item G.~H.~Golub and C.~F.~Van Loan, \emph{Matrix Computations}, 4th ed., Johns Hopkins, 2013, \S7.6, on why the Jordan
form is not computed in floating point; C.~Moler and C.~Van Loan, \emph{Nineteen dubious ways to compute the exponential of a
matrix, twenty-five years later}, SIAM Review 45 (2003), 3--49.
\item G.~Frobenius, \emph{\"Uber die mit einer Matrix vertauschbaren Matrizen}, Sitzungsberichte der Preussischen Akademie
(1910), 3--15.
\item N.~J.~Fine and I.~N.~Herstein, \emph{The probability that a matrix be nilpotent}, Illinois Journal of Mathematics 2
(1958), 499--504: there are exactly $q^{n(n-1)}$ nilpotent $n\times n$ matrices over $\mathbb F_q$ .
\end{itemize}
